\documentclass[a4paper,twoside]{article}


% --- RUIMTE EN LAYOUT ---
\usepackage{setspace}

\usepackage{geometry} %size regulating
\geometry{
 a4paper,
 left=25mm,
 right=25mm,
 top=20mm,
 bottom=20mm,
 }

\usepackage{parskip} % for better spacing between paragraphs
\usepackage{flushend} % Alleen nodig als je laatste pagina van artikel wilt vullen met tekst, zodat de kolommen even lang zijn. Anders kan je dit beter weglaten.

% --- WISKUNDE ENSYMBOLEN ---
\usepackage{amsmath, amssymb, amsfonts, amsthm} % For nicer math format
\usepackage{mathrsfs} % For curly arrr
\usepackage{bm} % For bold symbols eg \bm{\nabla}
\usepackage{esint} % For bolletje in integrals
\usepackage{dsfont}
\usepackage{siunitx} % Voeg deze expliciet toe voor die graden en eenheden!
\usepackage{braket} % For bra-ket notation
\usepackage{mathtools}
\usepackage{url}
\allowdisplaybreaks

% --- OVERIGE ---

\usepackage{graphicx} % Required for inserting images
\usepackage[hidelinks]{hyperref} % geen kaders om linken
\usepackage{subcaption}
%\usepackage[dutch]{babel}
\usepackage[T1]{fontenc}
\usepackage[utf8]{inputenc}
\usepackage{lmodern}
\usepackage{multicol}


% --- Tikz packages ---
\usepackage{tikz}
\usetikzlibrary{angles,quotes, babel} % for pic (angle labels)
\usetikzlibrary{arrows.meta,decorations.markings,calc, decorations.pathmorphing, decorations.pathreplacing, positioning, shapes, shadings, patterns}
\usepackage{xcolor}
\usepackage{tikz-3dplot}
\usepackage[siunitx]{circuitikz}

\usepackage{etoolbox} % ifthen
\usepackage[outline]{contour} % glow around text
\usepackage{xfp} % higher precision (16 digits?)
\contourlength{1.1pt}



% --- THEOREM STIJLEN MET EXTRA RUIMTE ---

\theoremstyle{definition}
\newtheorem{postulate}{Postulaat}
\newtheorem{definition}{Definitie}
\newtheorem{derivation}{Afleiding}
\newtheorem{proposition}{Propositie}
\newtheorem{example}{Voorbeeld}

\theoremstyle{plain}
\newtheorem{theorem}{Theorem}


% --- EXTRA COMMANDO'S ---

\newcommand{\dbar}{{\mkern+3mu\mathchar'26\mkern-12mu d}}

% --- Titel en Introductie --- 

\title{Kwantum II - Mijn oefeningen}
\author{T. Cools}
\date{2026-2027}

\begin{document}
\maketitle
\begin{abstract}
    In dit document hou ik bij hoe ik ervoor sta met de oefeningen van het vak Kwantum II.
\end{abstract}

\section{Formele ontwikkeling van de kwantummechanica}
\subsection{Oefening 1.1}
a) juist, maar misse fractie?\\
b) Juist\\
c) ?\\
\subsection{Oefening 1.2}
a) correct\\
b), c) en d) In de les aan bord gebracht.\\
\subsection{Oefening 1.3}
a) correct\\
b) geen idee, ik heb het niet zo uitgeeschreven\\
c) Heb ik wel gevonden\\
d) wss fout...
e) nog te doen
\subsection{Oefening 1.4}
Geen oplossingen van online?
\subsection{Oefening 1.5}
Geen oplossingen van online?
\subsection{Oefening 1.6}
Oplossing is yap
\subsection{Oefening 1.7}
Ik heb de eigenwaarden gevonden, maar mijn eigenvectoren zijn anders, ik heb die niet gecontrolleerd.
\subsection{Oefening 1.8}
Geen oplossing van online
\subsection{Oefening 1.9}
a) Hamiltoniaan in matrixvorm is juist.\\
b) Ik denk dat de oplossing \(  \hbar = 1\) heeft gekozen, dan heb ik dezelfde oplossing.\\
\subsection{Oefening 1.10}
Ziet er me correect uit!
\subsection{Oefening 1.11}
Geen oplossing van online
\subsection{Oefening 1.12} 
Gevonden, wel nog zelfstandig de waarden van \(A\) en \(B\) vinden in \(A\sin(\omega t) + B\cos(\omega t)\) met behulp van randvoorwaarden.     
\subsection{Oefening 1.13}
Geen oplossing van online.
\subsection{Oefening 1.14}
\subsubsection{Deelvraag a}
Mijn antwoord:
\[
\hat{p}(t) = \hat{p}(0)\cos(\omega t) - m \omega  \hat{x}(0) \sin(\omega t)
\]
\[
\hat{x}(t) = \hat{x}(0)\cos(\omega t) + \frac{\hat{p}(0)}{m \omega} \sin(\omega t)
\]

\textbf{Correct!}
\subsubsection{Deelvraag b}
\([\hat a , \hat{a}^\dagger] = \hat{\mathds{1}}\)
Had een rekenfouotje, zou normaal terug hetzelfde moeten uitkomen!
\subsection{Oefening 1.15}
Geen oplossing van online.
\subsection{Oefening 1.16}
Geen oplossing van online.
\section{Symmetrieën en behoudswetten in kwantummechanica}
\subsection{Oefening 2.1}
Ik kom uit op:
\[
[[\hat H , \hat x], \hat x] = -\frac{\hbar^2}{m} \hat{\mathds{1}}
\]
\[
[\hat H , \hat x] = -\frac{i\hbar \hat p}{m} \hat{\mathds{1}}
\]
Ik kom uit op:
\[
\sum_{n'} |\bra{\psi_{n}}\hat x \ket{\psi_{n'}}|^2 (E_{n'} - E_n) = \bra{\psi_n} [\hat x, \hat H \hat x] \ket{\psi_n}
\]
waar ik nu vast zit\dots
\subsection{Oefening 2.2}
Ik heb:
\[
\frac{d(\hat x \hat p)}{dt} = [\hat H, \hat x \hat p] = -\frac{i \hbar}{m} \hat p \hat p + i \hbar \hat x \bm \nabla (V(\hat x))
\]
\subsection{Oefening 2.3}   

\end{document}